برای بلوک اول (سیستم \lr{LTI}) داریم
$$w[n]=x[n]*h[n],\qquad h[n]\;\overset{\text{DTFT}}{\longleftrightarrow}\;H(\Omega)$$
و مطابق شکل:
$$
H(\Omega)=
\begin{cases}
1 & |\Omega|\le \frac{\pi}{2}\quad (\mathrm{mod}\ 2\pi)\\
0 & \text{سایر }
\end{cases}
$$
پس با خطی بودن،
$$x[n]=\delta[n]+1 \;\;\Longrightarrow\;\; w[n]=h[n]+1.$$

در بخش دوم دیاگرام:
$$y[n]=w[n]+(-1)^n w[n].$$
خروجی را به دو مؤلفه تجزیه می‌کنیم.

\begin{itemize}
\item مؤلفه‌ی ناشی از $\delta[n]$: \quad $w_\delta[n]=h[n]$ و
با خاصیت مدولاسیون $(-1)^n=e^{j\pi n}$ خواهیم داشت
$$(-1)^n h[n]\;\overset{e^{j\pi n}}{\longleftrightarrow}\;H(\Omega-\pi)$$
پس
$$Y_\delta(\Omega)=H(\Omega)+H(\Omega-\pi).$$
از آنجا که $H(\Omega)$ روی $[-\frac{\pi}{2},\frac{\pi}{2}]$ برابر ۱ است و شیفت $\pi$ آن بازه‌ی مکمل (در یک دوره) را پوشش می‌دهد، داریم
$$H(\Omega)+H(\Omega-\pi)=1 \quad \Longrightarrow\quad y_\delta[n]=\delta[n].$$

\item مؤلفه‌ی ناشی از ورودی ثابت $1$: \quad چون $H(0)=1$، خروجیِ بلوک اول برای ورودی DC برابر $w_1[n]=1$ است، لذا
$$y_1[n]=1+(-1)^n.$$
\end{itemize}

در نتیجه خروجی کل:
$$
y[n]=y_\delta[n]+y_1[n]=\delta[n]+1+(-1)^n.
$$